\section{Necessity and Order}
\label{sec:necessity}

The triad is a foundation only if each role carries independent load and if
later roles cannot safely manufacture their own earlier conditions. We apply
two conceptual tests: removal and order inversion.

\subsection{Removal tests}

\paragraph{Remove KFD-1.}
Suppose a system retains trust assessments and cooperation records but does
not preserve stable fact identities and change boundaries. A later participant
cannot determine whether an assessment still refers to the same contract,
artifact, occurrence, or meaning. Trust attaches to a moving object. Conflict
is resolved by the latest narrative or strongest authority because no stable
reference remains against which either can be challenged.

\paragraph{Remove KFD-2.}
Suppose facts are stable and participants are willing to cooperate, but no
purpose-bound assessment connects facts to reliance. Every participant must
either repeat the complete verification, accept another participant's
reputation, or act without knowing which gap matters. The archive may be
excellent, but reliance does not scale. Stable facts are necessary evidence;
they are not self-interpreting permission to act.

\paragraph{Remove KFD-3.}
Suppose facts and assessments are inspectable, but the system has no account of
participant value, choice, or visible constraint. It can support audit,
certification, and command execution. It cannot distinguish durable
cooperation from compliance obtained by monopoly, hidden workflow capture, or
unreviewable pressure. The world may be correct and trustworthy yet remain
non-generative across independent participants.

The failure modes are summarized in Table~\ref{tab:removal}.

\begin{table}[ht]
\centering
\small
\begin{tabular}{p{0.18\linewidth}p{0.31\linewidth}p{0.38\linewidth}}
\toprule
Removed role & What remains & Structural failure \\
\midrule
Stable facts & trust language and coordination & reliance binds to mutable
narrative or authority \\
Warranted trust & records and willing action & repeated verification, blind
reliance, or reputation substitution \\
Trusted-value cooperation & facts and assessment & accountable control may
remain, but independent joint generation collapses into compliance \\
\bottomrule
\end{tabular}
\caption{Removal tests for the three foundation roles.}
\label{tab:removal}
\end{table}

These tests establish functional non-redundancy, not mathematical uniqueness.
Another vocabulary may implement equivalent roles. The KFD claim is that the
capabilities themselves cannot be omitted from durable cooperation.

\subsection{Order-inversion failures}

The arrows in the triad describe dependency, not a rigid wall-clock sequence.
Participants may gather evidence, negotiate value, and act in overlapping
cycles. Nevertheless, a later layer may not retroactively create the basis on
which it claims to stand.

\paragraph{Trust before facts.}
When a system asks for reliance before exposing the relevant facts, identity,
reputation, confidence, or persuasive narrative becomes the substitute. Facts
may later be selected to defend the trust already granted. Correction then
requires defeating an authority relationship rather than updating an
assessment.

\paragraph{Cooperation before trustable value.}
When a system asks for commitment before value, facts, choices, and
constraints can be inspected, participation becomes the evidence that the
offer was acceptable. Lock-in is misread as preference and compliance as
agreement. The system can grow while losing the ability to learn why
participants would have chosen differently.

\paragraph{Successor facts admitted by the action that needs them.}
When completed action is allowed to rewrite its own prior fact cut, a result
can erase failed attempts, costs, reversals, or disconfirming evidence. The
loop becomes self-certifying: success defines which history counts. KFD-1
therefore applies again when new occurrences are admitted into the successor
world.

\subsection{Necessary, not sufficient}

The triad does not guarantee good outcomes. Facts may be incomplete, an
assessment mistaken, visible choices economically unequal, or cooperation
harmful to others outside the declared boundary. Domain knowledge, ethics,
law, security, governance, and empirical correction remain necessary. The
foundation claim is narrower: without continuity of fact, warranted reliance,
and participant-aware coordination, those additional mechanisms cannot form a
durable shared work world.
